%% =====================================================================
%%  CAPITULO 1 - INTRODUCAO
%% =====================================================================

\chapter{Introdução}
\label{cap:introducao}

Este capítulo apresenta o contexto do trabalho, o problema de engenharia
investigado, a justificativa, a hipótese, os objetivos, a delimitação do escopo,
uma visão resumida do método e a organização dos demais capítulos.

\section{Contextualização}
\label{sec:contextualizacao}

A produção de software assistida por modelos de linguagem de grande porte
(\textit{Large Language Models} --- LLM) deixou de ser experimento de laboratório
e passou a integrar rotinas reais de desenvolvimento. Modelos treinados em larga
escala demonstram capacidade de generalização para tarefas não vistas durante o
treinamento, fenômeno descrito como aprendizado em poucos exemplos
(\citeonline{brown2020}). A arquitetura de atenção
proposta por \citeonline{vaswani2017} viabilizou essa classe de modelos, e
técnicas de encadeamento de raciocínio ampliaram a capacidade de resolução de
problemas estruturados (\citeonline{wei2022}).

A partir dessa base, surgiram abordagens que organizam o modelo como agente
capaz de raciocinar e agir sobre um ambiente (\citeonline{yao2023react}), de
refletir sobre os próprios erros (\citeonline{shinn2023reflexion}) e de operar em
conjunto com outros agentes, cada um com papel especializado
(\citeonline{hong2024metagpt}; \citeonline{qian2024chatdev};
\citeonline{wu2023autogen}).

O ganho de velocidade é evidente; o ganho de qualidade é condicional. Um agente
que recebe um pedido vago produz código plausível, mas plausibilidade não é
corretude. Quando não existe especificação prévia, não há como decidir se o que
foi entregue é o que foi pedido; quando não existem critérios de aceitação, não
há como demonstrar que o pedido foi atendido; e quando não há registro de
evidência, a conclusão depende de confiança, não de verificação. A engenharia de
software já conhecia esse problema: a distinção entre verificação e validação e a
necessidade de critérios objetivos foram formalizadas há décadas
(\citeonline{boehm1984}), e a especificação de requisitos permanece sendo a
atividade com maior impacto sobre o custo final do produto
(\citeonline{sommerville2019}; \citeonline{pressman2016}).

\section{Problema de engenharia}
\label{sec:problema}

O problema investigado é:

\begin{citacao}
	Como estruturar um processo de desenvolvimento assistido por agentes de
	inteligência artificial de modo que cada entrega seja (i) rastreável a um
	requisito, (ii) verificável por um critério de aceitação e (iii) sustentada
	por evidência reproduzível?
\end{citacao}

O problema se decompõe em três dificuldades práticas. A primeira é de natureza
processual: agentes tendem a responder ao pedido mais recente, e não ao contrato
acumulado, de modo que sem um artefato de estado persistente o escopo se degrada
ao longo da execução. A segunda é de natureza epistemológica: a saída de um
agente é texto plausível, e plausibilidade não é evidência; é necessário
distinguir o que foi afirmado do que foi demonstrado. A terceira é de natureza
econômica: quanto maior o volume de contexto carregado, maior o custo e maior a
probabilidade de o agente tratar informação obsoleta como atual.

\section{Justificativa}
\label{sec:justificativa}

A relevância do estudo se sustenta em três argumentos.

\textbf{Argumento de risco.} As práticas que reduzem risco de regressão em
software --- integração contínua, testes automatizados, releases frequentes e
pequenas --- estão associadas a melhor desempenho organizacional
(\citeonline{humble2010}; \citeonline{forsgren2018}). Um processo que aceita
código gerado por agente sem passar pelos mesmos portões de qualidade importa
risco para dentro do produto.

\textbf{Argumento de auditabilidade.} Em domínios regulados ou críticos, a
pergunta relevante não é apenas se o software funciona, mas quem decidiu o quê,
com base em qual requisito e com qual evidência. A rastreabilidade é um requisito
de engenharia, não um subproduto (\citeonline{iso29148}).

\textbf{Argumento de reprodutibilidade.} Se duas execuções do mesmo pedido
produzem resultados diferentes e não existe especificação, o resultado não é
comparável nem acumulável. Sem artefatos congelados, o conhecimento produzido em
uma sessão se perde na seguinte.

Além disso, a literatura sobre agentes LLM aplicados à engenharia de software
concentra-se majoritariamente na geração de código
(\citeonline{hong2024metagpt}; \citeonline{qian2024chatdev}). O presente trabalho
desloca o foco da geração para o \textit{controle}: quais artefatos, portões e
condições de parada tornam a geração auditável.

\section{Hipótese}
\label{sec:hipotese}

A hipótese de trabalho é:

\begin{citacao}
	Um processo que obriga a especificação prévia com requisitos identificados,
	critérios de aceitação testáveis e evidência registrada produz entregas
	verificáveis por terceiros sem acesso ao histórico da conversa, ao custo de
	maior esforço documental.
\end{citacao}

A hipótese é falseável de duas maneiras. Primeiro, se algum critério de aceitação
não puder ser verificado sem consultar o autor, a rastreabilidade proposta
falhou. Segundo, se a evidência registrada não puder ser reproduzida a partir dos
artefatos versionados, o processo não é reprodutível. Ambos os testes foram
aplicados neste trabalho, e os resultados estão no Capítulo~\ref{cap:resultados}.

\section{Objetivos}
\label{sec:objetivos}

\subsection{Objetivo geral}
\label{subsec:objetivo-geral}

Definir, aplicar e avaliar um processo de desenvolvimento orientado por
especificações (\textit{spec-driven development} --- SDD) executado por agentes
de inteligência artificial, utilizando como estudo de caso a construção de um
jogo digital completo entregue em artefato único.

\subsection{Objetivos específicos}
\label{subsec:objetivos-especificos}

\begin{quadro}[htb]
	\caption{Objetivos específicos e critério de sucesso}
	\label{qua:objetivos}
	\centering
	\footnotesize
	\begin{tabularx}{\textwidth}{@{}p{1.1cm} X p{3.4cm}@{}}
		\toprule
		\textbf{ID} & \textbf{Objetivo específico} & \textbf{Critério de sucesso} \\
		\midrule
		OE-1 & Projetar o framework de agentes, instruções e artefatos de controle. & Framework descrito e versionado no repositório. \\
		\addlinespace
		OE-2 & Especificar formalmente o sistema-alvo com requisitos e critérios de aceitação verificáveis. & Todo requisito com identificador único. \\
		\addlinespace
		OE-3 & Registrar as decisões arquiteturais do alvo com alternativas descartadas. & Cada decisão com justificativa e alternativa rejeitada. \\
		\addlinespace
		OE-4 & Implementar o alvo sob as restrições especificadas. & Artefato único, sem dependências externas. \\
		\addlinespace
		OE-5 & Verificar cada critério de aceitação com evidência empírica. & Um resultado observável por critério. \\
		\addlinespace
		OE-6 & Avaliar o processo, suas limitações e ameaças à validade. & Ameaças classificadas e dados ausentes marcados. \\
		\bottomrule
	\end{tabularx}
	\legend{Fonte: elaborado pelo autor (2026).}
\end{quadro}

\section{Delimitação do escopo}
\label{sec:delimitacao}

\textbf{Dentro do escopo:} o processo de desenvolvimento, sua especificação,
sua aplicação a um sistema-alvo e a verificação empírica do alvo.

\textbf{Fora do escopo:} o treinamento ou a avaliação comparativa de modelos de
linguagem; a medição de custo financeiro por execução; a avaliação de usabilidade
com usuários humanos; e a comparação experimental contra grupo de controle, que
exigiria um segundo desenvolvimento paralelo do mesmo alvo.

Registra-se ainda uma delimitação factual relevante. O repositório onde o
trabalho foi conduzido contém a descrição de um sistema embarcado de controle
térmico, sem código, medição ou resultado correspondente. Como não há evidência
empírica desse sistema, ele não é objeto deste estudo; o estudo de caso é o
produto efetivamente implementado e verificado. Essa inconsistência documental
está registrada em \texttt{MONOGRAFIA\_PENDENCIAS.md} (P-26).

\section{Visão resumida do método}
\label{sec:metodo-resumido}

O trabalho é um estudo de caso único, de natureza aplicada, com abordagem
qualitativa predominante e componentes quantitativos pontuais
(\citeonline{yin2018}; \citeonline{runeson2009}). O procedimento teve quatro
etapas: (i) projeto do framework de agentes a partir das lacunas identificadas
em uma avaliação crítica prévia do material de origem; (ii) especificação do
alvo com requisitos funcionais, não funcionais e critérios de aceitação;
(iii) implementação sob restrições declaradas; e (iv) verificação empírica em
navegador, com inspeção estática e automação. O detalhamento está no
Capítulo~\ref{cap:metodos}.

\section{Organização do trabalho}
\label{sec:organizacao}

O Capítulo~\ref{cap:fundamentacao} apresenta a fundamentação teórica em
engenharia de software e agentes de inteligência artificial. O
Capítulo~\ref{cap:relacionados} compara trabalhos correlatos. O
Capítulo~\ref{cap:metodos} descreve materiais e métodos, incluindo ameaças à
validade. O Capítulo~\ref{cap:framework} detalha o framework SDD agêntico
proposto. O Capítulo~\ref{cap:especificacao} apresenta a especificação do
sistema-alvo. O Capítulo~\ref{cap:implementacao} descreve a implementação e suas
decisões arquiteturais. O Capítulo~\ref{cap:resultados} reporta experimentos e
resultados. O Capítulo~\ref{cap:discussao} discute os achados e as limitações. O
Capítulo~\ref{cap:conclusao} conclui e indica trabalhos futuros.
